\documentclass[12pt]{article}
\usepackage[margin=1in]{geometry}
\usepackage{enumitem}
\usepackage{titlesec}
\usepackage{parskip}
\usepackage{graphicx}
\usepackage{xcolor}
\usepackage{hyperref}
\usepackage{fontenc}

\title{\textbf{Jones and Hwang 2005 RG Notes}\\
\large Jones, Mark P., and Wonjae Hwang. ``Party Government in Presidential Democracies: Extending Cartel Theory Beyond the U.S. Congress'' [2005].\\
\textit{American Journal of Political Science} 49(2): 267-282.}
\author{}
\date{}

\begin{document}

\maketitle

In the United States, legislators have near-total control over their own careers; in much of the democratic world, this is not the case. Subparty bosses often have the opportunity to significantly influence or even entirely control legislators' careers (this is particularly true when elections depend on party lists).

Article is application of Cartel Theory to Argentinian politics (and, fourth-wall breaking, a commentary on applying US congressional studies to broader institutional questions).

Cartel Theory suggests that political parties have minimal impact on individual election results and subsequently have minimal leverage to control or influence legislators

Use Bayesianestimation to evaluate dimensionality of roll-call votes in the Argentinian legislature. Find very little influence of regional background/matters on voting behavior (beyond partisan influence/correlations)

overall: evidence of strong cartels in Argentinian legislature. Suggests applicability of cartel theory to certain parliamentary contexts and reinforces th pattern of power delegation to central party leaders in order to achieve coordinated policy outcomes

See Harmel \& Janda 1994 -- are parties democracy-seeking, and is this necessarily in opposition with cartel theory?

\end{document}
